\chapter{Resources and Further Reading}
\label{app:resources}

\section{Foundational papers}

\paragraph{CNN architectures.}
LeCun et al.\ (1998, LeNet-5) --- the conv-pool-FC template;
Krizhevsky et al.\ (2012, AlexNet) --- ReLU, dropout, GPU training;
Simonyan \& Zisserman (2014, VGG) --- depth from stacked $3\times3$;
He et al.\ (2015, ResNet) --- residual connections;
Tan \& Le (2019, EfficientNet) --- compound scaling;
Dosovitskiy et al.\ (2020, ViT) --- attention in vision.

\paragraph{Semantic segmentation.}
Long et al.\ (2014, FCN) --- first fully convolutional segmentation;
Ronneberger et al.\ (2015, U-Net) --- encoder--decoder with skips;
Chen et al.\ (2018, DeepLabV3+) --- ASPP + decoder;
Lin et al.\ (2017, FPN) --- feature pyramids;
Zhao et al.\ (2017, PSPNet) --- pyramid pooling.

\paragraph{EO / remote-sensing deep learning.}
Helber et al.\ (2019, EuroSAT); Wang et al.\ (2021, LoveDA);
Sumbul et al.\ (2019, BigEarthNet); Schmitt et al.\ (2019, SEN12MS);
Zhu et al.\ (2017) and Ma et al.\ (2019) surveys of DL in remote sensing;
Rolf et al.\ (2021) on distribution shift in satellite ML.

\paragraph{Transfer learning.}
Yosinski et al.\ (2014, transferability of features);
Neyshabur et al.\ (2020, what is transferred);
Marmanis et al.\ (2015, ImageNet$\to$EO transfer).

\section{Courses and documentation}

\begin{center}\small
\renewcommand{\arraystretch}{1.15}
\begin{tabular}{@{}ll@{}}
\toprule
Resource & Focus\\
\midrule
fast.ai Practical Deep Learning & practical DL with PyTorch\\
Stanford CS231n & CNN theory from fundamentals\\
PyTorch tutorials & framework implementation\\
torchgeo docs & EO datasets, GeoDatasets\\
SMP docs & segmentation model library\\
rasterio user guide & geospatial raster I/O\\
ESA EO College & remote-sensing background\\
\bottomrule
\end{tabular}
\end{center}

\section{Visualisation and GIS tools}

In-notebook: \texttt{matplotlib} (all plots), \texttt{seaborn} (confusion
matrices), \texttt{plotly}/\texttt{folium}/\texttt{ipyleaflet} (interactive
maps), \texttt{earthpy} (band composites), \texttt{rasterio.plot.show} (quick
GeoTIFF view). External GIS: QGIS (open prediction GeoTIFFs over OSM), Google
Earth Pro, ArcGIS Online, GeoServer (serve as WMS/WCS).

\section{Deployment ideas}

A FastAPI \texttt{/predict} endpoint that ingests a GeoTIFF and returns a
prediction map; a Streamlit dashboard (upload tile $\to$ map $\to$ download);
Google Earth Engine asset export; a Docker image based on a CUDA PyTorch base;
ONNX export for edge deployment on Jetson/OpenVINO.

\section{Portfolio suggestions}

Structure a public repo as \texttt{README} (banner + headline metric),
\texttt{model\_card.md} (architecture, training, limitations), executed
notebook, 2-page report, and an \texttt{assets/} folder (RGB tile, prediction
GeoTIFF, confusion matrix). A short GIF of the sliding-window inference
progressively covering a scene is a compelling LinkedIn/portfolio post.

\section{Future advanced topics}

Vision transformers for EO (SatMAE, Scale-MAE, Spectral-GPT);
self-supervised pretraining on Sentinel-2 (DINO, MAE);
change detection (FC-EF, BIT-CD); object detection in aerial imagery
(DOTA, oriented R-CNN, YOLO); time-series models (U-TAE, SITS-BERT);
EO foundation models (Prithvi, RemoteCLIP, GeoChat);
multi-modal SAR+optical and hyperspectral fusion; instance segmentation of
building footprints; continual learning across geographies.

\section{Dataset licenses}

\begin{center}\small
\renewcommand{\arraystretch}{1.15}
\begin{tabular}{@{}lll@{}}
\toprule
Dataset & License & Commercial use\\
\midrule
EuroSAT & MIT & yes\\
LoveDA & CC-BY-4.0 & yes (attribution)\\
BigEarthNet & CC-BY-4.0 & yes (attribution)\\
SpaceNet & CC-BY-SA-4.0 & attribution required\\
Sentinel-2 & Copernicus Open Access & yes\\
OpenStreetMap & ODbL & attribution required\\
Dynamic World labels & CC-BY-4.0 & yes (attribution)\\
\bottomrule
\end{tabular}
\end{center}

\vfill
\begin{center}\itshape
End of textbook. Now go and map the planet.
\end{center}
